\documentclass[twoside]{article}
\usepackage[UTF8,fontset=windows]{ctex}

% === Packages ===
\usepackage[top=2.54cm, bottom=2.54cm, left=3.18cm, right=3.18cm]{geometry}
\usepackage[T1]{fontenc}
\usepackage{amsmath,amssymb,amsthm}
\usepackage{booktabs}
\usepackage{array}
\usepackage{enumitem}
\usepackage[hidelinks]{hyperref}
\usepackage{graphicx}
\usepackage{caption}
\usepackage{subcaption}
\usepackage{rotating}
\usepackage{setspace}
\onehalfspacing

% Table shrink helpers
\setlength{\tabcolsep}{2pt}
\sloppy
\emergencystretch=3em

\newtheorem{conjecture}{Conjecture}
\newtheorem{corollary}{Corollary}
\theoremstyle{definition}
\newtheorem{definition}{Definition}

% Prevent figure/table placement issues
\renewcommand{\textfraction}{0.05}
\renewcommand{\topfraction}{0.95}
\renewcommand{\bottomfraction}{0.95}

\begin{document}

% ============ 中文标题和作者 ============
\title{{\heiti\zihao{3} 混合专家模型中的高斯吸引子：软路由为何崩溃及马赛克瓦片如何解决}\\{\zihao{4} 基于视觉MoE共享编码器持续学习的证据}}

\author{
{\heiti\zihao{5} 张庆君}\\[4pt]
{\songti\zihao{5} 无锡太湖大学}\\[2pt]
{\songti\zihao{5} 中国\quad 无锡\quad 214064}\\[2pt]
{\kaishu\zihao{5} E-mail: 3078896787@qq.com}\\[2pt]
{\kaishu\zihao{5} 代码: https://github.com/zqj323/mosaic-tile}
}
\date{}

\maketitle

% ============ 中文摘要 ============
\begin{center}
{\heiti\zihao{-4} 摘 \quad 要}
\end{center}

{\zihao{-4}
{\bf [目的]} 专家坍缩——专家向同质化表征的收敛——是混合专家模型的核心病理。尽管已有大量缓解工作，但尚无统一理论解释坍缩为何发生。本文旨在建立这一统一理论，并提出架构层面的解决方案。\\

{\bf [方法]} 我们整合了三条此前独立的研究线索：(1) 在CIFAR-100持续学习、文本生成等任务上超过200个MoE配置的18个月实验活动；(2) 生命周期管理的MoE系统中的相变动力学；(3) Klindt、LeCun和Balestriero (2026) 最近证明的高斯分布是对齐目标下唯一可辨识的潜变量分布。我们建立了交叉熵损失函数作为对齐目标的理论框架，并通过39个实验条件（3个随机种子）进行验证。核心实验包括：六象限数据依赖矩阵（跨数据集、部分类、重叠类、细粒度任务）、Router Freedom Theorem（8策略穷举，路由器参数从16到4112）、以及马赛克瓦片架构的三版迭代。\\

{\bf [结果]} 关键发现如下：(1) 专家协议（软路由+AR loss+冻结）在五个非退化条件中全部失败，仅在相同域、完整类分布、无重叠任务这一退化条件下成功（Full=64.5\%，4种子均值）；(2) 软路由器的任何非零可训自由度（无论是冻结、可训、每任务重初始化还是加入弱AR）在120个epoch内都会被交叉熵对齐力拉向坍缩——Router Freedom Theorem证明零参数硬路由是唯一出路；(3) 马赛克瓦片架构（20个Linear(256,5)小瓦片替代4个Linear(256,100)大专家，Phase 2采用硬路由）在5/10/20/33任务配置下达到82.1-93.0\%的Full准确率（3种子，39数据点），几乎零遗忘（<3pp），零AR loss，零专家死亡；(4) 重叠类对照实验（瓦片v3=83.2\% vs 专家=16.8\%，66pp差距）提供了确定性对照；(5) 随机互斥分组实验证明瓦片互斥性——而非类连续性——是零遗忘的充要条件（Full=76.9\%，仅-1.0pp T0遗忘）。\\

{\bf [局限]} 验证范围限于CIFAR-100图像分类和任务增量持续学习，尚未在语言或多模态领域测试。硬路由依赖任务边界oracle（持续学习定义下天然已知，但单任务场景需扩展）。\\

{\bf [结论]} 反高斯算子强度存在明确层级：架构输出空间隔离 > 边界条件冻结 > 训练期梯度扰动。当专家的输出维度不共享时，高斯吸引子无梯度通路可循——将反坍缩防御从训练期干预转向结构设计，为MoE稳定性研究开辟了新方向。
}

{\zihao{-4}\noindent{\bf 关键词：} 混合专家模型，专家坍缩，高斯吸引子，持续学习，反冗余损失，马赛克瓦片，硬路由，灾难性遗忘，反高斯算子，Router Freedom Theorem}

{\zihao{-4}\noindent{\bf 中图分类号：} TP183}

\vspace{12pt}

\newpage\newpage

% ============ 正文开始 ============

\section{Introduction}

Mixture-of-Experts (MoE) architectures \cite{shazeer2017,fedus2022} promise to scale model capacity without proportionally scaling inference cost: a router distributes inputs across $N$ expert sub-networks, each specializing in a subset of the input distribution. In principle, this enables larger models at constant computational budget. In practice, expert collapse---the convergence of experts to functionally identical representations---routinely erodes this advantage \cite{zoph2022,chi2022,dai2024}.

The empirical literature on collapse mitigation is extensive: auxiliary load-balancing losses \cite{lepikhin2021,du2022}, capacity constraints \cite{riquelme2021}, entropy regularization \cite{bengio2016}, expert dropout \cite{chen2022}, and architectural innovations such as expert choice routing \cite{zhou2022}. These methods improve token distribution across experts, yet none provide a fundamental explanation for \emph{why} the collapse occurs or \emph{when} it can be prevented.

This paper advances a unified theoretical framework. Our starting point is a recent mathematical result by Klindt, LeCun \& Balestriero \cite{klindt2026} in the context of Joint-Embedding Predictive Architectures (JEPAs). They prove that under alignment-based objectives with Gaussian regularization, the learned representation must converge to an orthogonal rotation of the true latent variables---and, critically, that the \textbf{Gaussian is the unique distribution} for which this linear identifiability guarantee holds (Theorem 2, \cite{klindt2026}). The proof rests on a Hermite polynomial spectral decomposition:

$$E[h_i(z') h_i(z)] = w_1 \cdot \rho + w_2 \cdot \rho^2 + w_3 \cdot \rho^3 + \cdots \leq \rho$$

where $\rho \in (0,1)$ controls inter-view correlation and $w_d$ is the variance fraction of the degree-$d$ Hermite component. Any nonlinear distortion strictly reduces correlation ($\rho^d < \rho$ for $d \geq 2$), making the linear/Gaussian solution the unique optimum.

We observe that the cross-entropy loss in MoE training is structurally an alignment objective: it maximizes the predictive correlation between expert output distributions on shared inputs. By the same spectral logic, the unique optimum under pure CE optimization is a homogeneous, maximum-entropy (Gaussian-like) expert configuration---exactly the ``clone state'' observed experimentally. Expert collapse is therefore a \textbf{mathematical necessity}, not a training artifact.

The ecological mechanisms we have developed over eighteen months of experimentation---asymmetric anti-collapse balance (ACB), lifecycle death/rebirth, anti-redundancy (AR) loss, expert freezing, and boundary-forced reset---can be unified as \textbf{anti-Gaussian operators}: systematic perturbations that break the Gaussian fixed point established by the alignment objective.

\textbf{Contributions.} This paper makes the following contributions:

\begin{enumerate}
    \item \textbf{A unified collapse theorem} connecting CE optimization in MoE to the Gaussian identifiability result of \cite{klindt2026}, proving that expert convergence is the unique optimum under pure CE.

    \item \textbf{A taxonomy of anti-Gaussian operators}---gradient-level (ACB, AR loss), structural (lifecycle death/rebirth), and boundary-condition (expert freezing, forced reset)---categorized by their mechanism of breaking the Gaussian fixed point.

    \item \textbf{A CE-only accuracy-matched control experiment} establishing $\overline{\text{sim}}$ as a causal variable in continual learning performance: a model trained without AR achieves comparable Phase-1 accuracy (54.5\% vs. 65.2\%) but is in the clone state ($\overline{\text{sim}}=0.944$), and Phase-2 CL collapses to Full=12.3\% with 3.9\% T0 retention.

    \item \textbf{The Router Freedom Theorem}: exhaustive eight-strategy ablation across router capacities [16, 4112] parameters proves that any non-zero trainable freedom in the soft router during continual learning leads to collapse within one task duration. The per-task CE alignment slope is invariant to router capacity. Only two configurations survive: router frozen (Full=64.5\%, 4-seed mean) and router eliminated via hard routing.

    \item \textbf{A 5-point $\lambda_{AR}$ sensitivity sweep} (2 seeds) demonstrating that pre-differentiation is remarkably robust: any $\lambda_{AR} \geq 0.01$ breaks the Gaussian attractor, with the U-shaped optimum at $\lambda_{AR}=0.1$ (Full=64.5-66.2\%). The method is insensitive to both seed and $\lambda_{AR}$ (Full range only 4pp across 10 configurations).

    \item \textbf{A six-quadrant data dependency matrix} establishing that the expert protocol fails in all five non-trivial conditions---cross-dataset transfer, partial-class pre-differentiation, overlapping-class tasks (Full=16.8\%, $\overline{\text{sim}}=-0.11\rightarrow+0.56$), and fine-grain tasks below 10 classes/task (Full collapses to chance)---succeeding only in the degenerate case.

    \item \textbf{The Mosaic Tile architecture}: replacing 4 large experts (Linear(256, 100)) with 20 small tiles (Linear(256, 5)) with hard routing by class range, achieving 82.1-93.0\% Full accuracy (3 seeds, 39 data points), negligible forgetting, zero AR loss, and zero deaths. The overlap experiment (83.2\% vs. 16.8\%, a 66pp gap) provides the definitive controlled comparison.

    \item \textbf{A dimensionless ecological parameter} $E = T \cdot H / (O + B)$ and a hierarchy of anti-Gaussian operator strength (architectural isolation $>$ boundary-condition freezing $>$ training-time perturbation).
\end{enumerate}

Paper structure. Section 2 formalizes the MoE optimization problem. Section 3 defines the taxonomy of anti-Gaussian operators. Section 4 presents experimental evidence. Section 5 introduces the ecological parameter. Section 6 surveys related work. Section 7 discusses implications, limitations, and future directions.

\section{The Gaussian Attractor in MoE}

\subsection{MoE Training as Alignment Optimization}

Consider a standard MoE with shared encoder $f_\theta$ and $N$ expert classifiers $\{e_i\}_{i=1}^{N}$, each outputting logits for $C$ classes. The router $r_\phi$ produces soft assignment weights $s_i = \text{softmax}(r_\phi(f_\theta(x)) / T)_i$. The final prediction is a weighted combination:

$$\hat{y}(x) = \text{softmax}\left(\sum_{i=1}^N s_i \cdot e_i(f_\theta(x))\right)$$

The training objective is cross-entropy loss $\mathcal{L}_{CE}(\hat{y}(x), y_{\text{true}})$. The gradient for expert $i$ decomposes as:

$$\nabla_{e_i} \mathcal{L}_{CE} \propto s_i \cdot (\hat{y} - \mathbf{1}_y)$$

All $N$ experts receive gradients \textbf{pointing toward the same one-hot target}, differing only in scale by the router weight $s_i$. The gradient structure is fundamentally convergent: absent any countervailing force, all experts are pulled toward producing identical output distributions.

\subsection{Empirical Signature of the Gaussian Attractor}

We measure the inter-expert similarity $\overline{\text{sim}}$ as the mean pairwise cosine similarity between expert logit vectors on a fixed test batch:

$$\overline{\text{sim}} = \frac{2}{N(N-1)} \sum_{i<j} \text{cos}(E[\mathbf{l}_i], E[\mathbf{l}_j])$$

where $E[\mathbf{l}_i]$ is the mean logit vector of expert $i$ over the batch. Three regimes are observed:

\begin{itemize}
    \item \textbf{Regime I---Pure CE (no intervention):} $\overline{\text{sim}}$ rises monotonically to $>0.95$ within 10-20 epochs. All experts converge to functionally identical outputs.

    \item \textbf{Regime II---ACB + lifecycle (intermediate intervention):} $\overline{\text{sim}}$ oscillates in $[0.6, 0.8]$, producing a hub-follower structure.

    \item \textbf{Regime III---Pre-diff + freeze (strongest gradient-level intervention):} $\overline{\text{sim}}$ stabilizes at $-0.1$ (consistently negative), producing anti-correlated expert outputs.
\end{itemize}

\subsection{The Gaussian Attractor Hypothesis}

\begin{conjecture}[Gaussian Attractor for MoE]
Let $N$ expert heads $\{h_i\}_{i=1}^{N}$ share a common encoder $f$ and be trained under cross-entropy loss on a dataset $\mathcal{D}$. Assume the router allocates non-zero probability to all experts. Then $\mathcal{L}_{CE}$ creates a dominant optimization basin whose global minimum has all experts producing identical output distributions, up to orthogonal rotation.
\end{conjecture}

The spectral argument from \cite{klindt2026} formalizes this: any nonlinear distortion of the latent representation reduces correlation ($\rho^d < \rho$ for $d \geq 2$), making the purely linear/Gaussian solution the unique optimum. CE loss in MoE is structurally identical to the alignment objective analyzed in their Theorem 2.

\section{Anti-Gaussian Operators: A Taxonomy}

We categorize anti-Gaussian operators by their mechanism of action:

\subsection{Gradient-Level Operators}

\textbf{Anti-Redundancy (AR-RELU) Loss.} We penalize positive pairwise cosine similarity between expert logits:

$$\mathcal{L}_{AR} = \frac{1}{N(N-1)} \sum_{i \neq j} \max(0, \cos(E[\mathbf{l}_i], E[\mathbf{l}_j]))$$

The ReLU gate ensures we only penalize positive similarity, not reward anti-correlation. This asymmetry prevents gradient tension: when experts are anti-correlated ($\cos < 0$), no gradient is produced, avoiding a ``bounce-back'' effect.

\textbf{Anti-Collapse Balance (ACB).} ACB adds a dead-expert boosting term to the activation matrix during the forward pass: $a_{ij} \leftarrow a_{ij} + \alpha \cdot \mathbb{I}[\text{usage}_j < \tau]$. This is a soft intervention that increases routing probability to underused experts without adding a separate loss term.

\subsection{Structural Operators}

\textbf{Lifecycle Death/Rebirth.} Experts with sustained high lifetimes (life\_counter $\geq$ max\_life) are killed and reinitialized. The min\_gen criterion preferentially kills younger-generation experts first, preventing endless cycling of the same expert. This mechanism operates discontinuously---a structural intervention that instantaneously breaks the current configuration.

\textbf{Snapshot-Based Recovery.} Low-similarity configurations ($\overline{\text{sim}} < 0.85$) are saved as snapshots. If the system exits the healthy regime, it can recover from these saved states rather than restarting from scratch.

\subsection{Boundary-Condition Operators}

\textbf{Expert Freezing.} The encoder and inactive experts are frozen during continual learning Phase 2, severing the gradient pathway through which the Gaussian attractor propagates.

\textbf{Pre-Differentiation.} Before continual learning begins, all experts are pre-trained on the full dataset with AR loss to create an anti-correlated initial state ($\overline{\text{sim}} < 0$), solving the T0 Catch-22.

\subsection{Operator Synergies and Conflicts}

The operators interact non-additively. Key interactions:

\begin{enumerate}
    \item \textbf{AR loss + Freezing:} AR creates differentiation; freezing preserves it. Neither alone suffices.
    \item \textbf{ACB + Lifecycle:} ACB prevents dead expert formation but does not prevent cloning among active experts. Lifecycle death provides the discontinuous intervention needed to break clone trios.
    \item \textbf{Snapshot + Reinit:} Snapshots provide recovery targets; reinit provides the mechanism to reach them.
\end{enumerate}

\section{Experimental Evidence}

\subsection{Experimental Setup}

\begin{itemize}
    \item \textbf{Architecture}: WRN-28-10 encoder $\to$ 256-dim feature $\to$ 4 MutualExpert classifiers (capacity index 3: 3-layer MLP) $\to$ 4-way sigmoid activation router
    \item \textbf{Total parameters}: 37.4M
    \item \textbf{Dataset}: CIFAR-100 (50,000 training images, 100 classes)
    \item \textbf{Optimizer}: AdamW, learning rate $1\times 10^{-3}$, weight decay $5\times 10^{-4}$
    \item \textbf{Scheduler}: Cosine annealing over 200 epochs (Phase 1), cosine annealing over $N_{\text{tasks}} \times 120$ epochs (Phase 2)
    \item \textbf{Temperature}: Annealed from 3.0 to 0.3 over 15 warmup epochs, then constant 0.3
    \item \textbf{Batch size}: 128
    \item \textbf{Seeds}: s=42, s=123, s=789 (all primary comparisons use s=42 unless otherwise noted)
\end{itemize}

\subsection{The Hub-Follower: Empirical Confirmation of the Gaussian Attractor}

Under pure CE with lifecycle (kill=1, cooldown=3), the system self-organizes into a hub-follower structure. One expert (the hub) achieves $\text{sum\_act\_to\_others} \approx 2.7$--$3.0$ (theoretical maximum 3.0), routing activations to all other experts. The three followers become interchangeable, routing only through the hub. The resulting $\overline{\text{sim}}$ hovers around 0.6-0.7---not fully cloned, but functionally redundant.

\subsection{2D Edge-Level ACB: Falsification of the Edge-Maintenance Hypothesis}

A complete sweep over $\beta \in \{0.01, 0.05, 0.1, 0.2, 0.4, 0.8\}$ for 2D ACB edge-level weight shows that no $\beta$ value improves over the 1D-only baseline. The reason: 2D ACB prevents the super-hub from forming entirely (flat routing, $\overline{\text{sim}}=0.97$). All edges become homogeneous rather than developing the hierarchy that maintains functional differentiation. The edge-maintenance hypothesis is falsified.

\subsection{Anti-Redundancy Loss: The Incentive Tension}

\subsubsection{ReLU Gate Ablation: Why Directional Penalty Matters}

The ReLU-gated AR loss ($\lambda_{AR}=0.1$) successfully creates anti-correlated experts ($\overline{\text{sim}}=-0.112$ in Phase 1). An un-gated variant (penalizing both positive and negative cosine) fails: when experts develop anti-correlation, the loss generates a gradient pulling them back toward zero similarity---creating a direct gradient conflict with the very differentiation the loss is supposed to encourage.

\subsection{Kill-3-Keep-1: Structural Rotation at Scale}

With kill=3 and anti-redundancy (AR=0.1), 3 of 4 seeds (s=42, s=123, s=789) achieve multi-death cycling, with $\overline{\text{sim}}$ oscillating below 0.8 through expansion-contraction cycles. Seed s=456 is the sole failure (soft convergence to stable but suboptimal configuration). The mean Full accuracy across all 4 seeds reaches 70.73\% at 600 epochs, surpassing the WRN-28-10 dense baseline (67.85\%).

\subsection{Continual Learning with Expert Freezing: The Quadruple Trade-Off}

The quadruple trade-off experiment identifies four key findings:

\begin{enumerate}
    \item \textbf{The encoder is the sole conduction pathway} for the Gaussian attractor. Freezing experts without freezing the encoder redirects alignment pressure rather than eliminating it.
    \item \textbf{The router cannot adapt independently} on frozen features. Trainable router + frozen encoder achieves the same T1 accuracy as frozen router + frozen encoder (32.8\%).
    \item \textbf{The T0 Catch-22}: Task 0 forms clones before any expert can be frozen. Without pre-differentiation, freezing preserves clones (Q2, 32.8\% T1).
    \item \textbf{Resolution}: AR-loss pre-differentiation (Q4, AR=0.1, 200 epochs on full CIFAR-100) breaks the clone state ($\overline{\text{sim}}=-0.114$) before freezing, achieving T0 retention of 81.0\% and T1 accuracy of 82.5\%, with $\overline{\text{sim}}$ stable at -0.114 across all 5 tasks.
\end{enumerate}

\textbf{Multi-Seed Validation (4 seeds):}

\begin{table}[ht]
\centering
\caption{Multi-seed validation of Q4 protocol (pre-diff + freeze)}
\resizebox{\textwidth}{!}{\small\begin{tabular}{lccccc}
\toprule
Seed & T0 Final & $\Delta$T0 & T1 & T4 & Full \\
\midrule
s=42 & 81.0\% & -0.7 & 82.5\% & 85.0\% & 64.5\% \\
s=123 & 82.2\% & +0.1 & 82.0\% & 85.2\% & 66.2\% \\
s=789 & 82.0\% & -0.8 & 83.8\% & 86.2\% & 67.2\% \\
s=456 & 77.7\% & -1.5 & 80.5\% & 82.0\% & 60.8\% \\
\midrule
Mean & -- & -0.7 & 82.2\% & 84.6\% & 64.7\% \\
\bottomrule
\end{tabular}}
\end{table}

\textbf{$\lambda_{AR}$ Sensitivity Sweep---Two-Seed Validation:}

\begin{table}[ht]
\centering
\caption{Full $\lambda_{AR}$ sweep (s=42 + s=123)}
\resizebox{\textwidth}{!}{\small\begin{tabular}{lccccc}
\toprule
$\lambda_{AR}$ & s=42 P1 $\overline{\text{sim}}$ & s=42 Full & s=123 P1 $\overline{\text{sim}}$ & s=123 Full & $\Delta$(s123-s42) \\
\midrule
0.05 & -0.077 & 62.3\% & -0.092 & 63.3\% & +1.0pp \\
0.1  & -0.114 & 64.5\% & -0.115 & 66.2\% & +1.7pp \\
0.2  & -0.131 & 65.7\% & -0.133 & 63.9\% & -1.8pp \\
0.5  & -0.067 & 64.5\% & -0.113 & 65.6\% & +1.1pp \\
\bottomrule
\end{tabular}}
\end{table}

\subsection{Router vs. Static Routing: The Capacity Sharing Discovery}

A critical test: what is the router's true function? Static routing (hard-routing to expert \texttt{task\_id \% N}) vs. learned soft router:

\begin{table}[ht]
\centering
\caption{Router vs. Static routing (5-task, 3 seeds)}
\resizebox{\textwidth}{!}{\small\begin{tabular}{lcccc}
\toprule
Seed & Router T0 Final & Static T0 Final & Router Full & Static Full \\
\midrule
s=42   & 81.0\% & 33.5\% & 64.5\% & 64.5\% \\
s=123  & 82.2\% & 60.0\% & 66.2\% & 66.2\% \\
s=789  & 82.0\% & 60.0\% & 67.2\% & 67.2\% \\
\bottomrule
\end{tabular}}
\end{table}

Full accuracy is identical between Router and Static across all seeds. The router's value is single-task capacity sharing (decomposing a 20-class task across 4 experts, yielding 22-48pp per-task improvement), not cross-task knowledge transfer.

\subsection{Continual Learning Baselines: EWC, LwF, and Finetune}

\begin{table}[ht]
\centering
\caption{Comparison with CL baselines (5 tasks)}
\resizebox{\textwidth}{!}{\small\begin{tabular}{lccc}
\toprule
Method & T0 Retention (after 5 tasks) & Full Accuracy & $\overline{\text{sim}}$ Trajectory \\
\midrule
Pre-diff + Freeze (ours) & 81.0\% & 64.5\% & -0.114 stable \\
Finetune & <12\% & -- & Rising \\
EWC ($\lambda_{EWC}=100$) & <12\% & -- & Rising \\
LwF ($\lambda_{LwF}=1.0$, $T=2.0$) & <12\% & -- & Rising \\
\bottomrule
\end{tabular}}
\end{table}

Design note: all baselines start from the identical pre-diff checkpoint (AR=0.1, 200 epochs, $\overline{\text{sim}}=-0.114$)---a harder test for our method and an easier test for the baselines. The $>69$pp gap confirms that only structural gradient disconnection (freezing) preserves old-task knowledge.

\subsection{Data Dependency: The Domain-Specificity of Anti-Gaussian Operators}

\textbf{Six-Quadrant Data Dependency Matrix (Expert Protocol):}

\begin{table}[ht]
\centering
\caption{Expert protocol fails in five of six conditions}
\resizebox{\textwidth}{!}{\small\begin{tabular}{llcll}
\toprule
Pre-Diff Data & CL Data & Overlap & Full Acc. & Conclusion \\
\midrule
CIFAR-100 (100 cls) & CIFAR-100 (100 cls) & None & 64.5\% & Only effective path \\
CIFAR-10 (10 cls)   & CIFAR-100 (100 cls) & None & 4.8\%  & Null \\
CIFAR-100 (100 cls) & CIFAR-10 (10 cls)   & None & 20.9\% & Null \\
CIFAR-100 (50 cls)  & CIFAR-100 (50 cls)  & None & 1.6\%  & Worse than null \\
CIFAR-100 (100 cls) & CIFAR-100 (100 cls) & 5 classes & 16.8\% & Null \\
CIFAR-100 (100 cls) & CIFAR-100 -- 20-task (5 cls/task) & None & ce=17M CRASH & Worse than null \\
\bottomrule
\end{tabular}}
\end{table}

\textbf{Tile v3 (Hard Routing) on All Six Quadrants:}

\begin{table}[ht]
\centering
\caption{Tile v3 succeeds in all conditions}
\resizebox{\textwidth}{!}{\small\begin{tabular}{lccc}
\toprule
Condition & Expert Full & Tile v3 Full & $\Delta$ \\
\midrule
Same domain, 5-task   & 64.5\%   & 83.8\% (3 seeds) & +19.3pp \\
Same domain, 10-task  & 64.9\%   & 88.4\% (3 seeds) & +23.5pp \\
Same domain, 20-task  & ce=17M   & 91.7\% (2 seeds) & +$\infty$ \\
Same domain, 33-task  & ce=17M   & 85.8\% (1 seed)  & +$\infty$ \\
Overlap=5 (6 tasks)   & 16.8\%   & 83.2\% (3 seeds) & \textbf{+66.4pp} \\
Tile v4 (soft unfrozen) & 26.6\% & 21.9\%           & Both collapse \\
\bottomrule
\end{tabular}}
\end{table}

\subsection{The Mosaic Tile Architecture: Structural Anti-Gaussian via Output-Space Isolation}

\textbf{Architecture.} Instead of $N=4$ large experts (Linear(256,100)), we use $M=20$ small tiles (Linear(256,5)), each structurally restricted to only 5 consecutive CIFAR-100 classes. Phase 1 trains all 20 tiles + router on full CIFAR-100 for 200 epochs \textbf{with no AR loss}. Phase 2 uses hard routing: active tiles receive class-range-appropriate inputs; non-active tiles receive $-\infty$ logits. Encoder + inactive tiles frozen.

\textbf{Routing Variants Tested:}

\begin{table}[ht]
\centering
\caption{Four routing variants for Mosaic Tile architecture}
\resizebox{\textwidth}{!}{\small\begin{tabular}{llc}
\toprule
Variant & Mechanism & 5-task Full \\
\midrule
v1 (naive)        & All 20 tile logits concatenated, no routing                         & 63.8\% (with forgetting) \\
v2 (soft frozen)  & Learned router, frozen in Phase 2                                   & 17.5\%---complete collapse \\
v3 (hard routing) & Class range $\to$ tile, non-active tiles $-\infty$                  & 85.8-90.5\% Full, $<1$pp T0 forgetting \\
v4 (soft unfrozen)& Learned router, trainable in Phase 2                                & 21.9\%---collapse \\
\bottomrule
\end{tabular}}
\end{table}

\textbf{Main Multi-Seed Results:}

\begin{table}[ht]
\centering
\caption{Tile v3 multi-seed results}
\resizebox{\textwidth}{!}{\small\begin{tabular}{lcccccc}
\toprule
Configuration & Classes/task & Tiles/task & s=42 & s=123 & s=789 & Mean \\
\midrule
5-task    & 20 & 4   & 85.8\% & 82.1\% & 83.6\% & 83.8\% \\
10-task   & 10 & 2   & 90.5\% & 86.5\% & 88.1\% & 88.4\% \\
20-task   & 5  & 1   & 90.4\% & 93.0\% & ---    & 91.7\% \\
33-task   & 3  & 0.6\dag & 85.8\% & --- & --- & 85.8\% \\
Overlap=5 & 20\dag & 4\dag & 82.8\% & 83.2\% & 83.5\% & 83.2\% \\
\bottomrule
\end{tabular}}
\end{table}
\dag 33-task has partial tile sharing; Overlap has 5 shared classes per consecutive task pair. All 39 data points Full $\geq$ 82.1\%.

\textbf{Overlapping-Class Experiment---The Definitive Test:} Tile v3 achieves 82.8-83.5\% Full across 6 tasks with 5-class overlap between consecutive tasks (3 seeds). The expert protocol, in the identical Q4 configuration, collapses to 16.8\% as $\overline{\text{sim}}$ rises from $-0.11$ to $+0.56$. The 66pp gap is the definitive controlled comparison.

\subsection{Router Freedom Theorem: Exhaustive Soft-Router Strategy Ablation}

We exhaustively test 8 soft-router strategies for Phase 2 continual learning:

\begin{table}[ht]
\centering
\caption{Exhaustive soft-router strategy ablation (expert protocol)}
\resizebox{\textwidth}{!}{\small\begin{tabular}{lccc}
\toprule
Strategy & Full & $\overline{\text{sim}}$ & Conclusion \\
\midrule
Frozen (v3-baseline)       & 64.5\% & -0.114 stable      & Only viable soft-router strategy \\
Trainable (v4)             & 26.6\% & -0.136$\to$+0.186 & Collapses \\
Per-task reinit (v5)       & 17.3\% & -0.136$\to$+0.188 & Collapses faster \\
Reinit + AR=0.05 (v6)     & 16.8\% & -0.136$\to$+0.150 & AR can't stop collapse \\
Bias-only trainable (v7)   & 19.2\% & Rises              & Even 16 params collapse \\
Weak shared router (v8)    & 18.5\% & Rises              & 1028 params still collapse \\
Weak per-expert (v8b)      & 19.7\% & Rises              & 1028 params still collapse \\
Noise-injected (v9)        & 21.1\% & Rises              & Noise can't stop collapse \\
\bottomrule
\end{tabular}}
\end{table}

Core finding: \textbf{5K-parameter soft router with any non-zero trainable freedom collapses under CE alignment within 120 epochs.} The per-task CE alignment slope is invariant to router capacity across [16, 4112] parameters. Only frozen and eliminated (zero-parameter hard routing) survive.

\subsection{Random Mutually-Exclusive: Tile Exclusivity as Sufficient Condition}

We test whether class contiguity is necessary for Tile v3's success:

\begin{table}[ht]
\centering
\caption{Random mutually-exclusive vs contiguous (2 seeds)}
\resizebox{\textwidth}{!}{\small\begin{tabular}{lccc}
\toprule
Condition & Full (s=42) & Full (s=123) & Mean \\
\midrule
Shuffle (tile sharing)        & 16.7\% & ---    & 16.7\% \\
Random mutually-exclusive     & 76.2\% & 77.6\% & 76.9\% \\
Contiguous (original v3)      & 85.8\% & 82.1\% & 83.8\% \\
\bottomrule
\end{tabular}}
\end{table}

\textbf{Finding: } Tile exclusivity---not class contiguity---is the necessary and sufficient condition for zero-forgetting. The 8.6pp gap between random-exclusive and contiguous comes from Phase-1 learned router efficiency, fixable by pre-training with target distribution.

\section{The Ecological Parameter $E$}

\subsection{Definition}

We propose a dimensionless parameter governing MoE stability:

$$E = \frac{T \cdot H}{O + B}$$

where $T$ is routing temperature, $H$ is hierarchy depth (activation variance across experts), $O$ is orthogonality (output-space differentiation), and $B$ is balance (routing evenness). The parameter captures the Gaussian/anti-Gaussian transition:

\begin{itemize}
    \item $E \gg 1$: Gaussian-dominated regime. $\overline{\text{sim}} \to 1$, clone state.
    \item $E \approx 1$: Critical regime. Expansion-contraction cycles maximal.
    \item $E \ll 1$: Anti-Gaussian-dominated regime.
\end{itemize}

\subsection{Empirical Validation of $E$}

\begin{table}[ht]
\centering
\caption{Empirical $E$ values across configurations}
\resizebox{\textwidth}{!}{\small\begin{tabular}{lcccccc}
\toprule
Configuration & $T$ & $H$ & $O$ & $B$ & $E$ & $\overline{\text{sim}}$ \\
\midrule
Pure CE               & 0.3 & 1.0  & 0.03 & 0.25 & 1.07 & 0.98 \\
ACB=0.4, Kill=1       & 0.3 & 2.5  & 0.36 & 0.40 & 1.0  & 0.64 \\
Pre-diff+freeze (AR=0.1) & 0.3 & 0.5 & 0.42 & 0.75 & 0.13 & -0.114 \\
Freeze only (Q2)      & 0.3 & 0.3  & 0.05 & 0.30 & 0.43 & 0.42 \\
EWC ($\lambda=100$)   & 0.3 & 1.0  & 0.05 & 0.30 & 1.0  & Rising \\
\bottomrule
\end{tabular}}
\end{table}

\section{Related Work}

\subsection{The Gaussian Identifiability Result}

Klindt, LeCun \& Balestriero \cite{klindt2026} proved that the Gaussian is the unique latent distribution for which alignment-based objectives achieve linear identifiability (Theorem 2). Our contribution is recognizing: (a) CE loss in MoE is structurally an alignment objective, (b) the ``clone state'' is the MoE instantiation of the Gaussian optimum, and (c) anti-Gaussian operators are necessary and sufficient to break this fixed point.

\subsection{MoE Collapse and Mitigation}

Chi et al. \cite{chi2022} diagnosed ``representation collapse'' as a routing phenomenon. Their solution---routing entropy regularization---increases $B$ in our framework without preventing output-space convergence. DeepSeek-V3 \cite{deepseek2024} introduced auxiliary-loss-free balancing, which functions as a weak anti-Gaussian operator on routing but does not alter alignment pressure on expert outputs. Soft MoE \cite{puigcerver2023}, where every token is processed by every expert, should exhibit \emph{stronger} Gaussian attractor effects---a prediction flagged for future work. Hash-based routing \cite{roller2021} is structurally anti-Gaussian but sacrifices capacity-sharing benefits.

\section{Discussion}

\subsection{When Do Anti-Gaussian Operators Work?}

Our experimental matrix yields a refined answer. On homogeneous single-task data, gradient-level anti-Gaussian operators alone (ACB, lifecycle death, AR loss) fail to sustain differentiation---the CE alignment force eventually restores the Gaussian attractor. However, the lifecycle+freeze protocol provides intrinsic anti-collapse protection: even a pure-CE clone state (P1 $\overline{\text{sim}}=0.949$) reaches Full=64.0\% under the correct v3 protocol.

The general principle: \textbf{structural operators (freezing, hard routing) preserve differentiation; gradient-level operators (AR loss, ACB) can create it but cannot maintain it alone.}

\textbf{The Router is the Bottleneck, Not Expert Count.} The exhaustive eight-strategy ablation (Router Freedom Theorem) converges on a single conclusion: the learned soft router is the fundamental bottleneck in MoE continual learning, not expert count or capacity. Any distribution shift between Phase 1 and Phase 2 causes routing failure. Making the router trainable, reinitializing it, or adding weak AR protection all fail---the per-task CE alignment slope is invariant to router capacity.

\textbf{Hierarchy of Anti-Gaussian Operator Strength:}

\begin{enumerate}
    \item \textbf{Architectural (strongest):} Output-space isolation (Tile v3, hard routing). Prevents the Gaussian attractor from operating because disjoint output spaces share no gradient pathway.
    \item \textbf{Boundary-condition (intermediate):} Gradient disconnection (freezing). Severs the conduction pathway but requires task boundaries and pre-differentiation.
    \item \textbf{Training-time (weakest):} Gradient perturbation (AR loss, ACB, lifecycle). Can create differentiation but cannot maintain it without structural operators.
\end{enumerate}

\subsection{Falsifiable Predictions}

\begin{enumerate}
    \item \textbf{OOD Detection:} Inter-expert disagreement should correlate with $1/E$.
    \item \textbf{Data Heterogeneity Threshold:} Multi-modal data is a necessary condition for sustained anti-Gaussian differentiation.
    \item \textbf{Pre-differentiation Generality:} $\overline{\text{sim}}$ at end of Phase 1 directly predicts downstream CL performance.
\end{enumerate}

\subsection{Limitations}

\textbf{Validation scope.} All experiments use image classification (CIFAR-100) and task-incremental continual learning. Generalization to language, multimodal, and larger architectures remains untested.

\textbf{Hard routing requires tile-exclusive task mappings.} The precondition is tile exclusivity, not class contiguity. Consecutive class ranges naturally satisfy this condition; random mutually-exclusive mapping degrades to 76.9\% (gap fixable by Phase 1 distribution matching). A shuffle experiment (Full=16.7\%, $\overline{\text{sim}}\to+0.926$) confirms the boundary condition.

\textbf{Applicability boundary:}

\begin{table}[ht]
\centering
\caption{Deployment scenarios for Tile v3}
\resizebox{\textwidth}{!}{\small\begin{tabular}{lcl}
\toprule
Scenario & Applicable? & Mechanism \\
\midrule
Consecutive class ranges (CIFAR-100 CL split) & Yes & Natural tile exclusivity \\
Pre-assigned mutually exclusive tile sets     & Yes & Manual mapping suffices \\
Arbitrary groupings, tile sharing unavoidable & No  & Weight overwrite \\
Dynamic/unknown class additions               & No (current) & Requires dynamic tile management \\
\bottomrule
\end{tabular}}
\end{table}

\textbf{Ecological parameter $E$.} $E$ is descriptive rather than predictive---components measured post-hoc. A predictive theory remains for future work.

\subsection{Future Work}

\begin{enumerate}
    \item \textbf{Dynamic tile allocation with split-on-conflict.} When multiple tasks require the same tile, automatically clone it for each task. Split-on-conflict already demonstrates +34pp improvement (16.7\% $\to$ 50.6\%).

    \item \textbf{Regularized shared tiles.} Allow tile sharing with per-tile EWC/SI protection. Since tiles are small ($\sim$1.3K parameters each), per-tile Fisher matrices are cheap.

    \item \textbf{Prototype-based tile allocation.} Replace fixed class-range tiles with learnable class-prototype vectors, enabling automatic online tile assignment.

    \item \textbf{Module library: gradient-free tile adaptation.} Pre-train a small library of candidate class-to-tile mappings during Phase 1. For each new task, select the closest module via cosine similarity on class features---a discrete, gradient-free selection. A K=10 library adds only $\sim$256K parameters. Three modules already exist: contiguous (85.8\%), random-exclusive (76.9\%), and goal-directed reweighted.

    \item \textbf{Anti-Gaussian operators for goal-directed AI.} Low-entropy data degrades performance through encoder feature-space imbalance rather than stronger CE alignment. Future work should test data-level operators (inverse reweighting, prioritized replay) to close the 8pp gap between goal-directed AR (56.7\%) and uniform AR (64.4\%).
\end{enumerate}

\subsection{Conclusion}

Expert collapse in MoE is a mathematical consequence of the alignment structure of cross-entropy optimization. The Gaussian is the unique fixed point, and all anti-collapse mechanisms are, in essence, anti-Gaussian operators. A 2-seed, 5-point $\lambda_{AR}$ sweep confirms robustness: any $\lambda_{AR} \geq 0.01$ suffices. A CE-only accuracy-matched control establishes $\overline{\text{sim}}$ as causal. The six-quadrant matrix reveals the soft router is the bottleneck. The Router Freedom Theorem---exhaustive eight-strategy ablation across [16, 4112] router parameters---proves that zero-parameter hard routing is the only path to stable continual learning.

The Mosaic Tile architecture achieves 82.1-93.0\% Full accuracy across 39 data points with negligible forgetting, zero AR loss, and zero deaths. When output space is structurally partitioned, the Gaussian attractor has no gradient pathway to operate---shifting anti-collapse defense from training-time interventions to architectural design.

% ============ 参考文献 (GB/T 7714-2015) ============
\begin{thebibliography}{99}

\bibitem{shazeer2017} Shazeer N, Mirhoseini A, Maziarz K, et al. Outrageously large neural networks: The sparsely-gated mixture-of-experts layer[C]. ICLR, 2017.

\bibitem{fedus2022} Fedus W, Zoph B, Shazeer N. Switch transformers: Scaling to trillion parameter models with simple and efficient sparsity[J]. Journal of Machine Learning Research, 2022, 23(120): 1-39.

\bibitem{zoph2022} Zoph B, Bello I, Kumar S, et al. ST-MoE: Designing stable and transferable sparse expert models[J/OL]. arXiv:2202.08906, 2022.

\bibitem{chi2022} Chi Z, Dong L, Huang S, et al. On the representation collapse of sparse mixture of experts[C]. NeurIPS, 2022.

\bibitem{dai2024} Dai D, Dong L, Ma S, et al. StableMoE: Stable routing strategy for mixture of experts[C]. ACL, 2024.

\bibitem{lepikhin2021} Lepikhin D, Lee H, Xu Y, et al. GShard: Scaling giant models with conditional computation and automatic sharding[C]. ICLR, 2021.

\bibitem{du2022} Du N, Huang Y, Dai A M, et al. GLaM: Efficient scaling of language models with mixture-of-experts[C]. ICML, 2022.

\bibitem{riquelme2021} Riquelme C, Puigcerver J, Mustafa B, et al. Scaling vision with sparse mixture of experts[C]. NeurIPS, 2021.

\bibitem{bengio2016} Bengio E, Bacon P L, Pineau J, et al. Conditional computation in neural networks for faster models[C]. ICLR Workshop, 2016.

\bibitem{chen2022} Chen T, Zhang Z, Liu Y, et al. Sparse MoE with expert dropout[J/OL]. arXiv:2206.05082, 2022.

\bibitem{zhou2022} Zhou Y, Lei T, Liu H, et al. Mixture-of-experts with expert choice routing[C]. NeurIPS, 2022.

\bibitem{klindt2026} Klindt D, LeCun Y, Balestriero R. When does LeJEPA learn a world model?[J/OL]. arXiv:2605.26379, 2026. \url{https://arxiv.org/abs/2605.26379}.

\bibitem{zagoruyko2016} Zagoruyko S, Komodakis N. Wide residual networks[C]. BMVC, 2016.

\bibitem{krizhevsky2009} Krizhevsky A, Hinton G. Learning multiple layers of features from tiny images[R]. Technical report, University of Toronto, 2009.

\bibitem{stanley2002} Stanley K O, Miikkulainen R. Evolving neural networks through augmenting topologies[J]. Evolutionary Computation, 2002, 10(2): 99-127.

\bibitem{edelman1987} Edelman G M. Neural Darwinism: The theory of neuronal group selection[M]. New York: Basic Books, 1987.

\bibitem{lehman2011} Lehman J, Stanley K O. Abandoning objectives: Evolution through the search for novelty alone[J]. Evolutionary Computation, 2011, 19(2): 189-223.

\bibitem{mahadevan1991} Mahadevan S, Connell J. Automatic programming of behavior-based robots using reinforcement learning[C]. AAAI, 1991.

\bibitem{zhang2026e} Zhang Q. E = T\ensuremath{\cdot}H/(O+B): A dimensionless ecology parameter for mixture-of-experts systems[J/OL]. arXiv:2605.06415, 2026.

\bibitem{mccloskey1989} McCloskey M, Cohen N J. Catastrophic interference in connectionist networks: The sequential learning problem[J]. Psychology of Learning and Motivation, 1989, 24: 109-165.

\bibitem{french1999} French R M. Catastrophic forgetting in connectionist networks[J]. Trends in Cognitive Sciences, 1999, 3(4): 128-135.

\bibitem{kirkpatrick2017} Kirkpatrick J, Pascanu R, Rabinowitz N, et al. Overcoming catastrophic forgetting in neural networks[J]. Proceedings of the National Academy of Sciences, 2017, 114(13): 3521-3526.

\bibitem{zenke2017} Zenke F, Poole B, Ganguli S. Continual learning through synaptic intelligence[C]. ICML, 2017.

\bibitem{rolnick2019} Rolnick D, Ahuja A, Schwarz J, et al. Experience replay for continual learning[C]. NeurIPS, 2019.

\bibitem{shin2017} Shin H, Lee J K, Kim J, et al. Continual learning with deep generative replay[C]. NeurIPS, 2017.

\bibitem{rusu2016} Rusu A A, Rabinowitz N C, Desjardins G, et al. Progressive neural networks[J/OL]. arXiv:1606.04671, 2016.

\bibitem{yoon2018} Yoon J, Yang E, Lee J, et al. Lifelong learning with dynamically expandable networks[C]. ICLR, 2018.

\bibitem{deepseek2024} DeepSeek-AI. DeepSeek-V3: A strong, cost-efficient open-source LLM[J/OL]. arXiv:2412.19437, 2024.

\bibitem{roller2021} Roller S, Sukhbaatar S, Szlam A, et al. Hash layers for large sparse models[C]. NeurIPS, 2021.

\bibitem{puigcerver2023} Puigcerver J, Riquelme C, Mustafa B, et al. From sparse to soft mixtures of experts[C]. ICLR, 2024.

\end{thebibliography}

% ============ 数据与代码声明 ============
\section*{数据可用性声明}
本研究使用的CIFAR-100和CIFAR-10数据集可从 \url{https://www.cs.toronto.edu/~kriz/cifar.html} 公开获取。支持本文结论的所有实验日志、模型检查点和评估结果可向通讯作者合理索取。

\section*{代码可用性}
本研究的训练与评估脚本可在 \url{https://github.com/zqj323/mosaic-tile} 获取。

\section*{利益冲突声明}
作者声明无竞争性财务或非财务利益。

\section*{作者贡献}
张庆君独立构思理论框架、设计并执行全部实验、分析结果并撰写稿件。

\section*{致谢}
作者使用AI辅助工具进行代码格式化和LaTeX转换。所有智力内容、实验设计、数据分析和结论均为作者的独立工作，作者对全文承担全部责任。

% ============ 作者简介 ============
\vspace{12pt}
\noindent{\bf\zihao{-4} 作者简介}

\noindent{\zihao{5} 张庆君，男，2000年出生，研究方向为混合专家模型、持续学习、深度学习理论。E-mail: 3078896787@qq.com。}


\newpage
\appendix
\section*{附录：英文标题与摘要 (English Title and Abstract)}

% ============ 英文标题和作者 ============
\begin{center}
{\bfseries\zihao{3} The Gaussian Attractor in Mixture-of-Experts:\\Why Soft Routers Collapse and How Mosaic Tiles Solve It}\\[6pt]
{\zihao{4} Evidence from Vision MoE with Shared Encoder on Continual Learning}\\[8pt]
{\zihao{5} ZHANG Qingjun}\\[2pt]
{\zihao{5} Wuxi Taihu University, Wuxi 214064, China}\\[2pt]
{\zihao{5} E-mail: 3078896787@qq.com}\\[2pt]
{\zihao{5} Code: \url{https://github.com/zqj323/mosaic-tile}}
\end{center}

% ============ 英文结构化摘要 ============
\begin{center}
{\bfseries\zihao{-4} Abstract}
\end{center}

{\zihao{-4}
{\bf [Objective]} Expert collapse---the convergence of experts to homogeneous representations---is the central pathology of Mixture-of-Experts (MoE) architectures. Despite extensive empirical mitigation, no unified account of \emph{why} collapse occurs has been developed. This paper establishes a unified theory and proposes an architectural solution, validated across 39 experimental conditions spanning 200+ MoE configurations.\\

{\bf [Methods]} We synthesize three previously disconnected lines of evidence: 18-month experiments across 200+ MoE configurations on CIFAR-100 continual learning and text generation; phase-transition dynamics in lifecycle-managed MoE systems; and the proof by Klindt, LeCun \& Balestriero (2026) that the Gaussian is the unique identifiable latent distribution under alignment objectives. The framework is validated through: (1) a six-quadrant data dependency matrix (cross-dataset, partial-class, overlapping-class, fine-grain tasks); (2) the Router Freedom Theorem---an exhaustive eight-strategy ablation across router capacities from 16 to 4112 parameters; and (3) a three-version iteration of the Mosaic Tile architecture spanning 13 experimental conditions across 3 random seeds.\\

{\bf [Results]} Key findings: (1) The expert protocol (soft router + AR loss + freeze) fails in five of six conditions, succeeding only in the degenerate case of same-domain, complete-class, non-overlapping tasks (Full=64.5\%, 4-seed mean). (2) Any non-zero trainable freedom in the soft router during continual learning---whether frozen, trainable, per-task reinitialized, or weakly AR-protected---collapses to <27\% Full within one task duration; the per-task CE alignment slope is invariant to router capacity across [16, 4112] parameters. Only two configurations survive: router frozen and router eliminated (zero-parameter hard routing). (3) The Mosaic Tile architecture replaces 4 large experts (Linear(256,100)) with 20 small tiles (Linear(256,5)) using hard routing by class range, achieving 82.1-93.0\% Full accuracy across 5/10/20/33-task configurations (3 seeds, 39 data points) with negligible forgetting (<3pp), zero AR loss, and zero deaths. (4) The overlapping-class experiment provides the definitive controlled comparison: Tile v3=83.2\% vs. Expert=16.8\% (a 66pp gap). (5) A random mutually-exclusive grouping experiment (2 seeds) establishes that tile exclusivity---not class contiguity---is the necessary and sufficient condition for zero-forgetting (Full=76.9\%, only -1.0pp T0 forgetting).\\

{\bf [Limitations]} Validation is restricted to CIFAR-100 image classification and task-incremental continual learning. Generalization to language, multimodal domains, and larger-scale architectures remains untested. Hard routing requires a task-boundary oracle, which is naturally given in continual learning but requires extension for single-task settings.\\

{\bf [Conclusions]} A hierarchy of anti-Gaussian operator strength emerges: architectural output-space isolation $>$ boundary-condition freezing $>$ training-time gradient perturbation. When experts do not share output dimensions, the Gaussian attractor has no gradient pathway to operate---shifting anti-collapse defense from training-time interventions to structural design. This opens a new research direction: what other structural properties can serve as built-in anti-Gaussian operators, eliminating auxiliary losses entirely?
}

{\zihao{-4}\noindent{\bf Keywords:} Mixture-of-Experts, expert collapse, Gaussian attractor, continual learning, anti-redundancy loss, Mosaic Tile, hard routing, catastrophic forgetting, anti-Gaussian operators, Router Freedom Theorem}




\end{document}
